%%%%%%%%%%%%%%%%%%%%%%%%%%%%%%%%%%%%%%%%%%%%%%%%%%%%%%%%%%%%%%%%%%%%%%%%%%%%%
%% 03_arch_clock_concept.tex
%% Section: Clock Concept
%% Chapter: Architecture Concepts
%%
%% Describes the Clock Generation Unit (CGU) shared by both chip variants
%% Variant B (RV_PipelineCacheFwd); clock requirements not yet defined.
%%
%% Integration in 00_rwurv64i_all.tex:
%%   \section{Clock Concept}
%%   \input{\TEXTarch 03_arch_history_tab04}
%%   \input{\TEXTarch 03_arch_clock_concept}
%%
%% Requires packages: tabularx, booktabs, float
%%
%% v0.1  2026-07-12  Claude-AI  Initial version -- RV_NoPipeline clock requirements
%% v0.2  2026-11-26  Claude-AI  CGU register interface on clk_div_s
%%%%%%%%%%%%%%%%%%%%%%%%%%%%%%%%%%%%%%%%%%%%%%%%%%%%%%%%%%%%%%%%%%%%%%%%%%%%%

%%%%%%%%%%%%%%%%%%%%%%%%%%%%%%%%%%%%%%%%%%%%%%%%%%%%%%%%%%%%%%%%%%%%%%%%%%%%%
\subsection{Clock Generation Unit (CGU) Overview}
\label{subsec:cguOverview}
%%%%%%%%%%%%%%%%%%%%%%%%%%%%%%%%%%%%%%%%%%%%%%%%%%%%%%%%%%%%%%%%%%%%%%%%%%%%%

Both chip variants (Table~\ref{tab:asChipVariants}) share the same Clock
Generation Unit IP (\texttt{as\_cgu\_top} / \texttt{as\_cgucore}).
The chip has a single external clock input, \texttt{clk\_i}, from which the
CGU derives all internal clock domains using free-running counter-based
dividers. Every divided clock leaves the CGU directly from a flip-flop (the
registered compare \texttt{cnt < div/2}): a decoded counter value could
glitch while the counter bits switch, and a glitch on a clock line is an
extra clock edge. The registered output has the same waveform, delayed by one
\texttt{clk\_i} cycle.
There is no PLL: every derived clock is an integer division of \texttt{clk\_i}.

\begin{table}[H]
\caption{CGU Clock Domains}
\label{tab:asCguDomains}
\centering
\begin{tabularx}{\textwidth}{|l |l |X |X|}
  \hline
  Output & Divider (fixed) & Consumers & Notes \\
  \hline
  \hline
  \texttt{clk\_core\_o} & \texttt{clk\_core\_div} = 80 &
    CPU core, GPIO0, UART0, ICU0, Scratchpad, D-/I-Cache, CGU register interface &
    Functional clock for the CPU and all Wishbone slaves
    (as \texttt{clk\_div\_s}, Section~\ref{subsubsec:cguRegClock}) \\
  \hline
  \texttt{clk\_bus1\_o} & \texttt{clk\_bus1\_div} = 80 & Reserved & Not currently consumed in \texttt{as\_top\_mem} \\
  \hline
  \texttt{clk\_bus2\_o} & \texttt{clk\_bus2\_div} = 100 & Reserved & Not currently consumed in \texttt{as\_top\_mem} \\
  \hline
  \texttt{clk\_qspi\_o} & \texttt{clk\_qspi\_div} = 4 & QSPI controller & \\
  \hline
  \texttt{tck\_i} & --- (external JTAG clock) & JTAG TAP, boundary scan &
    Fully asynchronous to \texttt{clk\_i}; no timing relationship is analysed
    between the two domains \\
  \hline
\end{tabularx}
\end{table}

\textbf{The divider ratios are fixed at RTL compile time, not run-time
configurable.}
Unlike the QSPI and UART peripherals (which expose a genuine \texttt{CLKDIV}
Wishbone register), the CGU implements only a read/write peripheral
\textbf{ID register}; the four divider ratios above are SystemVerilog
\texttt{localparam}s in \texttt{as\_pack.sv} and take effect unconditionally.
Changing a divider therefore requires an RTL/constant change and a full
re-synthesis -- it cannot be adjusted by software or by writing a register.

\subsubsection{Scan-Test Clock Bypass}

During normal operation the CPU and Wishbone peripherals are clocked by
\texttt{clk\_div\_s}, which is generated in \texttt{as\_top\_mem} as:
\begin{verbatim}
assign clk_div_s = dr_cap_s ? clk_i : clk_core_s;
\end{verbatim}
\texttt{dr\_cap\_s} is driven by the JTAG TAP.
In normal function mode (\texttt{dr\_cap\_s}\,=\,0) the CPU runs on the
divided \texttt{clk\_core\_s}.
In scan-test mode (\texttt{dr\_cap\_s}\,=\,1) the divider is bypassed and the
CPU is clocked directly by the raw external \texttt{clk\_i}, so that an ATE
tester can single-step or free-run the scan chain independently of the CGU
divide ratio.
The CGU dividers (\texttt{as\_cgucore}) always run on the undivided
\texttt{clk\_i}; only their divided outputs are subject to the bypass mux.

\subsubsection{CGU Register Clock}
\label{subsubsec:cguRegClock}

\texttt{as\_cgu\_top} has two clock inputs: \texttt{clk\_i} drives only
the divider counters, \texttt{clk\_wb\_i} the slave BPI and the ID
register. Both variants connect \texttt{clk\_wb\_i} to
\texttt{clk\_div\_s}, so the CGU register interface is in the clock
domain of the CPU and the other Wishbone slaves. Up to 2026-11-25 the ID
register was clocked by \texttt{clk\_i}; a CPU write then crossed into
the \texttt{clk\_i} domain (found by the first XH018 STA of Variant~A).

%%%%%%%%%%%%%%%%%%%%%%%%%%%%%%%%%%%%%%%%%%%%%%%%%%%%%%%%%%%%%%%%%%%%%%%%%%%%%
\subsection{Clock Requirements -- Variant B}
\label{subsec:clockReqB}
%%%%%%%%%%%%%%%%%%%%%%%%%%%%%%%%%%%%%%%%%%%%%%%%%%%%%%%%%%%%%%%%%%%%%%%%%%%%%

Not yet defined. The testbenches run at the same operating point as
Variant~A (\texttt{clk\_i}\,=\,100\,MHz, \texttt{clk\_core\_s}\,=\,1.25\,MHz).
